% commands.tex -- mathematical shorthand for this paper

% Theorem environments (numbered within sections)
\newtheorem{theorem}{Theorem}[section]
\newtheorem{lemma}[theorem]{Lemma}
\newtheorem{corollary}[theorem]{Corollary}
\newtheorem{proposition}[theorem]{Proposition}

% Definitions and remarks (unnumbered or section-numbered as you prefer)
\theoremstyle{definition}
\newtheorem{definition}[theorem]{Definition}
\theoremstyle{remark}
\newtheorem{remark}[theorem]{Remark}

% TODO marker macro -- visible red [TODO: ...] in draft, suppress for release.
% Defined with \long so the argument may contain \par / itemize / multi-
% paragraph prose.  Uses {\color{...} ...} grouping rather than \textcolor
% because the latter is not \long.
\makeatletter
\long\def\todomark#1{{\color{red}[TODO: #1]}}
\long\def\placeholder#1{\par\noindent\todomark{#1}\par}
\makeatother
% To suppress all TODO markers for a release build, uncomment:
% \long\def\todomark#1{}
% \long\def\placeholder#1{}

% Add framework-specific shorthand below.  Keep names short, semantic, and
% consistent across sections.  Examples:
%   \newcommand{\R}{\mathbb{R}}                          % a math symbol
%   \newcommand{\op}[1]{\hat{#1}}                        % a wrapper
%   \newcommand{\eqdef}{\stackrel{\mathrm{def}}{=}}      % a relation
